\hspace{3mm}

\centerline{\bf \Huge Вступительный экзамен по алгебре}
\centerline{\bf \Huge }

\begin{flushright}
        \scriptsize{\textbf{Математика — это единственный совершенный метод водить самого себя за нос. Альберт Эйнштейн}}
\end{flushright}

\problem{\textbf{1.1}} Решите систему из двух линейных алгебраических уравнений(СЛАУ):

\begin{equation*}
 \begin{cases}
   17x - y = 13
   \\
   9y - 153x = -118
 \end{cases}
\end{equation*}

\begin{flushright} \textbf{3 балла}\\ 
\end{flushright}

\problem{\textbf{1.2}} Решите систему из двух линейных алгебраических уравнений(СЛАУ):

\begin{equation*}
 \begin{cases}
   2x + y = 1
   \\
   3x - 7y = -24
 \end{cases}
\end{equation*}

\begin{flushright} \textbf{3 балла}\\ 
\end{flushright}

\problem{\textbf{1.3}} Решите систему из двух линейных алгебраических уравнений(СЛАУ):

\begin{equation*}
 \begin{cases}
   -3x + 7y = 13
   \\
   57x - 133y = -247
 \end{cases}
\end{equation*}

\begin{flushright} \textbf{3 балла}\\ 
\end{flushright}

\problem{\textbf{2.1}} Чему равно следующее выражение: 

$40.08\cdot17.56^{2}-308.3536\cdot17.56+13.36^{3}\ -13.36^{2}\cdot52.68$

\begin{flushright} \textbf{6 баллов}\\ 
\end{flushright}

\problem{\textbf{2.2}} Чему равно следующее выражение: 

$\dfrac{20\dfrac{8}{15}\cdot7.5-54.6:\dfrac{2}{5}}{3\dfrac{13}{21}\cdot8.4-34.4:14\dfrac{1}{3}}+43.75:11\dfrac{2}{3}+24.6:1\dfrac{1}{5}$

\begin{flushright} \textbf{6 баллов}\\ 
\end{flushright}

\hspace{3mm}

\problem{\textbf{3.1}} Решите систему из двух линейных алгебраических уравнений(СЛАУ):

\begin{equation*}
 \begin{cases}
   \dfrac{1}{13}x + 2.2y = -3\dfrac{1}{12}
   \\
   4.1x - 5\dfrac{1}{3}y = 6.17
 \end{cases}
\end{equation*}

\begin{flushright} \textbf{6 баллов}\\ 
\end{flushright}

\hspace{3mm}

\problem{\textbf{3.2}} Решите систему из трёх линейных алгебраических уравнений(СЛАУ):

\begin{equation*}
 \begin{cases}
   x + 2y + 3z = 12
   \\
   2x - y + z = -1
   \\
   3x + 7y - 2z = 5
   
 \end{cases}
\end{equation*}

\begin{flushright} \textbf{6 баллов}\\ 
\end{flushright}

\hspace{3mm}

\problem{\textbf{4}} Прямая $Q$ проходит через точки $A(-7\dfrac{8}{3};11)$ и $B(4;-2.5)$.

Найдите уравнение, прямой проходящей через точку $C(-7;7)$, перпендикулярной прямой $Q$.

\hspace{3mm}

\begin{flushright} \textbf{6 баллов}\\ 
\end{flushright}

\hspace{3mm}

\problem{\textbf{5}} АА купил своим ученикам 20 кг сникерсов и 30 кг киндеров, заплатив за всё 210 тысяч рублей. Если сникерсы подорожают на 30\%, а киндеры подешевеют на 10\%, то за такую же массу соответствующих сладостей АА заплатит 213 тысяч рублей. Сколько рублей стоит 100г каждой сладости?


\begin{flushright} \textbf{6 баллов}\\ 
\end{flushright}

\hspace{3mm}

\hspace{3mm}

\problem{\textbf{6}} Найдите пересечение и объединение пяти множеств:

\begin{flushleft}
  $1) 11 \leqslant x   < 100$
\end{flushleft}

\begin{flushleft}
  $2) -333 < x  < 75$
\end{flushleft}

\begin{flushleft}
  $3) 10 < x  \leqslant 99$
\end{flushleft}

\begin{flushleft}
  $4) -456 < x  < -333$
\end{flushleft}

\begin{flushleft}
  $5) -332 \leqslant x  \leqslant 74$
\end{flushleft}

\begin{flushright} \textbf{6 баллов}\\ 
\end{flushright}

\problem{\textbf{7}} Решите задачу на движение:

Два пешехода вышли на рассвете. Каждый шёл с постоянной скоростью. Один шёл из $A$ в $B$, другой - из $B$ в $A$. Они встретились в полдень и, не прекращая движения, пришли: один в $B$ в 4 часа вечера, а другой - в $A$ в 9 часов вечера. В котором часу в тот день был рассвет?

\begin{flushright} \textbf{9 баллов}\\ 
\end{flushright}

\problem{\textbf{8}} Решите задачу на логику:

На острове ЭлМШ живут рыцари и лжецы. Рыцари всегда говорят правду, лжецы всегда лгут. Некоторые жители заявили, что на острове четное число рыцарей, а остальные заявили, что на острове нечетное число лжецов. Может ли число жителей острова быть нечетным?

\begin{flushright} \textbf{9 баллов}\\ 
\end{flushright}









